\documentclass{handout}

\assignment{Homework 3}
\title{Circular motion}
\duedate{Monday September 22, 2025}

\begin{document}
\maketitle

\hwinstructions

\begin{questions}
    \question \textbf{Race track speed limit:} There is only so much force tires can provide from friction before they start to skid. Knowing this, you hope to do a public service and put up a speed limit sign at your local race track so nobody skids. You know the following information about the cars:
    \begin{itemize}
        \item The cars are all heavier than 1,500 kg.
        \item The cars' tires can all provide 23 kN of friction force before skidding.
        \item The radius of the turns on the track are 150 m.
    \end{itemize}
    \begin{parts}
        \part What number should you put on the speed limit sign? (your answer will probably be in m/s but try converting it to miles per hour).
        \part Draw the three forces on the car if it is on \textbf{flat} ground. You should have friction, normal force, and gravity. Circle the force that counts as \textbf{centripetal}.
        \part Draw the three forces on the car if it is on a \textbf{banked} turn. Which forces now contribute to the \textbf{centripetal} force?
        \part Explain why having a banked turn would let you raise the speed limit on the track.
    \end{parts}

    \question \textbf{Loop-de-loop}: You are building a marble drop track and you want to add a loop-de-loop. You're pretty sure that you can get the marble going 1.5 m/s at the top of the loop. We need to figure out how big the loop should be.
    \begin{parts}
        \part Calculate the centripetal acceleration needed to go around a 20 cm radius loop.
        \part Calculate the centripetal acceleration needed to go around a 25 cm radius loop.
        \part Calculate the radius of the motion if the centripetal acceleration all coming from the gravitational force on the marble.
        \part Draw two pictures: a track with radius 20 cm and a track with 25 cm radius. On both, draw a dotted line representing the circular motion which would occur under only the influence of gravity.
        \part Should you build the 20 cm or the 30 cm loop? Which forces will be providing the \textbf{centripetal force} on the marble at the top of the loop in the case of your choice?
    \end{parts}

    \question \textbf{Periods}: Suppose we know a mass is moving in a circular motion, and it is being kept in that motion by the tension from a string.
    \begin{parts}
        \part If the string is 1 m long and you are spinning the mass around 4 times a second. What is the force provided by the string if the object has a mass of 50 g?
        \part Suppose your answer to part (a) was 1/4 of the breaking force of the string. What is the smallest circle you could spin the object in \textbf{with the same period} before the string breaks?
    \end{parts}
\end{questions}

\end{document}
